% Generated by make_companion_labels.sh from schubert_bialgebra_separate_paper.aux on 2026-10-15T18:34Z. Do not edit.
\begin{filecontents*}[overwrite]{schubert_bialgebra_separate_paper_labels.aux}
\relax 
\providecommand\hyper@newdestlabel[2]{}
\providecommand\HyField@AuxAddToFields[1]{}
\providecommand\HyField@AuxAddToCoFields[2]{}
\citation{ressayre2009branching}
\@writefile{toc}{\contentsline {section}{\tocsection {}{1}{Introduction}}{1}{section.1}\protected@file@percent }
\@writefile{toc}{\contentsline {subsection}{\tocsubsection {}{}{The Schubert bialgebra and its dual}}{1}{section*.1}\protected@file@percent }
\citation{postnikov2009chains}
\citation{nadeau2024quasisymmetric}
\citation{bergeron2025qsymflag}
\citation{notes}
\citation{little2003combinatorial}
\citation{ressayre2009branching}
\@writefile{toc}{\contentsline {subsection}{\tocsubsection {}{}{Littlewood--Richardson rules for dual bases}}{2}{section*.2}\protected@file@percent }
\newlabel{theorem:LR}{{1.1}{2}{LR rule for dual Schubert polynomials}{theorem.1.1}{}}
\citation{bergeron2025qsymflag}
\citation{mathieu1990filtrations}
\citation{assaf2017slide}
\citation{samuelforest}
\citation{lam2021back}
\citation{Yu2025EmbeddingBP}
\citation{Weigandt2021}
\citation{GaoHuang2023}
\citation{HuangShimozonoYu2024}
\@writefile{toc}{\contentsline {subsection}{\tocsubsection {}{}{Relation to the forest Littlewood--Richardson rule}}{3}{section*.3}\protected@file@percent }
\@writefile{toc}{\contentsline {subsection}{\tocsubsection {}{}{Further directions}}{3}{section*.4}\protected@file@percent }
\citation{assaf2018demazure}
\citation{nadeau2024forest}
\@writefile{toc}{\contentsline {subsection}{\tocsubsection {}{}{Outline of the paper}}{4}{section*.5}\protected@file@percent }
\citation{sottile}
\citation{sottile}
\citation{notes}
\citation{samuelleibniz}
\@writefile{toc}{\contentsline {section}{\tocsection {}{2}{The Schubert Bialgebra and its dual}}{5}{section.2}\protected@file@percent }
\newlabel{section:bialgebra}{{2}{5}{The Schubert Bialgebra and its dual}{section.2}{}}
\@writefile{toc}{\contentsline {subsection}{\tocsubsection {}{2.1}{Some preliminaries}}{5}{subsection.2.1}\protected@file@percent }
\newlabel{definition:pieri}{{2.1.1}{5}{Two Pieri relations $\tom {k}$ and $\Tom {k}$}{lemma.2.1.1}{}}
\newlabel{item:thintom}{{4}{5}{Two Pieri relations $\tom {k}$ and $\Tom {k}$}{Item.4}{}}
\newlabel{item:varthintom}{{2.1.1}{5}{Two Pieri relations $\tom {k}$ and $\Tom {k}$}{Item.4}{}}
\newlabel{definition:divdiff}{{2.1.2}{5}{Divided differences and skew divided differences}{lemma.2.1.2}{}}
\citation{bsskew}
\citation{coprod}
\citation{bsskew}
\newlabel{definition:codestuff}{{2.1.3}{6}{Lehmer code, dual code, dominant permutations}{lemma.2.1.3}{}}
\newlabel{definition:polysequence}{{2.1.4}{6}{Pulling out variables from double Schubert polynomials, the relation $\downvar {i}$}{lemma.2.1.4}{}}
\newlabel{proposition:pullindex}{{2.1.6}{7}{}{lemma.2.1.6}{}}
\@writefile{toc}{\contentsline {subsection}{\tocsubsection {}{2.2}{Definition}}{7}{subsection.2.2}\protected@file@percent }
\newlabel{remark:fatidentity}{{2.2}{7}{Fat identities}{subsection.2.2}{}}
\@writefile{toc}{\contentsline {subsection}{\tocsubsection {}{2.3}{The Schubert basis}}{8}{subsection.2.3}\protected@file@percent }
\citation{sottile}
\newlabel{section:elementary}{{2.4}{9}{The elementary basis}{subsection.2.4}{}}
\@writefile{toc}{\contentsline {subsection}{\tocsubsection {}{2.4}{The elementary basis}}{9}{subsection.2.4}\protected@file@percent }
\newlabel{lemma:powerdiff}{{2.4.3}{9}{}{lemma.2.4.3}{}}
\citation{notes}
\newlabel{lemma:doublediff}{{2.4.5}{10}{}{lemma.2.4.5}{}}
\newlabel{definition:elemcoeff}{{2.4.6}{10}{}{lemma.2.4.6}{}}
\newlabel{theorem:elemtrans}{{2.2}{10}{Transition to the elementary basis}{theorem.2.2}{}}
\citation{schubmult}
\@writefile{toc}{\contentsline {section}{\tocsection {}{3}{The dual algebra}}{11}{section.3}\protected@file@percent }
\newlabel{section:dualalgebra}{{3}{11}{The dual algebra}{section.3}{}}
\citation{notes}
\@writefile{toc}{\contentsline {subsection}{\tocsubsection {}{3.1}{The dual Schubert basis}}{12}{subsection.3.1}\protected@file@percent }
\newlabel{definition:dualpath}{{3.1.1}{12}{}{lemma.3.1.1}{}}
\newlabel{theorem:dualformula}{{3.1}{12}{}{theorem.3.1}{}}
\newlabel{lemma:dpieri}{{3.1.3}{12}{}{lemma.3.1.3}{}}
\newlabel{eq:branching}{{3.1.1}{13}{The dual Schubert basis}{equation.3.1.1}{}}
\newlabel{lemma:dualbranch}{{3.1.4}{13}{}{lemma.3.1.4}{}}
\@writefile{toc}{\contentsline {subsection}{\tocsubsection {}{3.2}{The coproduct}}{13}{subsection.3.2}\protected@file@percent }
\newlabel{theorem:dschcoproduct}{{3.2}{13}{Schubert LR coefficients as dual Schubert coproduct}{theorem.3.2}{}}
\newlabel{example:dschcoproduct}{{3.2.1}{13}{}{lemma.3.2.1}{}}
\citation{bbrc}
\@writefile{toc}{\contentsline {section}{\tocsection {}{4}{The ring of bounded RC graphs}}{14}{section.4}\protected@file@percent }
\newlabel{section:brc}{{4}{14}{The ring of bounded RC graphs}{section.4}{}}
\@writefile{toc}{\contentsline {subsection}{\tocsubsection {}{4.1}{The module of bounded RC graphs}}{14}{subsection.4.1}\protected@file@percent }
\newlabel{definition:rcgraph}{{4.1.1}{14}{}{lemma.4.1.1}{}}
\newlabel{lemma:d1desc}{{4.1.2}{14}{\cite [Corollary~3.11]{bbrc}}{lemma.4.1.2}{}}
\newlabel{lemma:factor}{{4.1.3}{14}{}{lemma.4.1.3}{}}
\newlabel{theorem:modactrc}{{4.1}{15}{}{theorem.4.1}{}}
\newlabel{lemma:drcpieri}{{4.1.5}{15}{}{lemma.4.1.5}{}}
\newlabel{corollary:evschub}{{4.1.6}{16}{}{lemma.4.1.6}{}}
\@writefile{toc}{\contentsline {subsection}{\tocsubsection {}{4.2}{The pipe dream visualization and roots}}{16}{subsection.4.2}\protected@file@percent }
\citation{little2003combinatorial}
\citation{billey2019bijective}
\citation{GaoHuang2023}
\@writefile{lof}{\contentsline {figure}{\numberline {1}{\ignorespaces The pipe dream visualization of the bounded RC graph $R$ with $\mathtt  {ht}(R)=5$ where $R=\{(1,1),(1,2),(2,1),(3,1),(3,3)\}$}}{17}{figure.caption.6}\protected@file@percent }
\providecommand*\caption@xref[2]{\@setref\relax\@undefined{#1}}
\newlabel{figure:rcexamppipe}{{1}{17}{The pipe dream visualization of the bounded RC graph $R$ with $\hr (R)=5$ where $R=\{(1,1),(1,2),(2,1),(3,1),(3,3)\}$}{figure.caption.6}{}}
\@writefile{lof}{\contentsline {figure}{\numberline {2}{\ignorespaces An invalid set of crossings causing pipes to cross more than once, caused by inserting a crossing at a negative root}}{17}{figure.caption.7}\protected@file@percent }
\newlabel{section:bpd}{{4.3}{17}{Zeroing out the last row}{subsection.4.3}{}}
\@writefile{toc}{\contentsline {subsection}{\tocsubsection {}{4.3}{Zeroing out the last row}}{17}{subsection.4.3}\protected@file@percent }
\citation{lam2021back}
\citation{lam2021back}
\citation{GaoHuang2023}
\citation{GaoHuang2023}
\citation{little2003combinatorial}
\citation{billey2019bijective}
\citation{huang2020monk}
\citation{GaoHuang2023}
\citation{GaoHuang2023}
\@writefile{lof}{\contentsline {figure}{\numberline {3}{\ignorespaces The map $\mathcal  {Z}$ on $R=\{(1,2),(1,3),(2,2)\}$, of height $3$ and permutation $1432$. The crossing $(2,2)$ moves to $(2,1)$; the resulting word $3\,2\,2$ is not reduced, so $(1,2)$ moves to $(1,1)$, giving the reduced word $3\,1\,2$ of $1432\,t_{34}t_{13}=2413$. The empty row $3$ is then deleted.}}{18}{figure.caption.8}\protected@file@percent }
\newlabel{fig:zeroing}{{3}{18}{The map $\zeromap $ on $R=\{(1,2),(1,3),(2,2)\}$, of height $3$ and permutation $1432$. The crossing $(2,2)$ moves to $(2,1)$; the resulting word $3\,2\,2$ is not reduced, so $(1,2)$ moves to $(1,1)$, giving the reduced word $3\,1\,2$ of $1432\,t_{34}t_{13}=2413$. The empty row $3$ is then deleted}{figure.caption.8}{}}
\@writefile{toc}{\contentsline {subsubsection}{\tocsubsubsection {}{}{Bumpless pipe dreams}}{18}{section*.9}\protected@file@percent }
\newlabel{definition:littlebump}{{4.3.1}{18}{Little bumps}{lemma.4.3.1}{}}
\@writefile{toc}{\contentsline {subsubsection}{\tocsubsubsection {}{}{Monk moves}}{18}{section*.10}\protected@file@percent }
\newlabel{eq:monk}{{4.3.1}{18}{Monk moves}{equation.4.3.1}{}}
\citation{GaoHuang2023}
\citation{little2003combinatorial}
\citation{huang2020monk}
\newlabel{eq:GHintertwine}{{4.3.2}{19}{Monk moves}{equation.4.3.2}{}}
\newlabel{lemma:bpdconfine}{{4.3.2}{19}{Row confinement}{lemma.4.3.2}{}}
\newlabel{lemma:bumpconfine}{{4.3.3}{19}{Little bumps and top rows}{lemma.4.3.3}{}}
\newlabel{lemma:bpdcompletion}{{4.3.4}{19}{Bounded completion}{lemma.4.3.4}{}}
\newlabel{proposition:toprows}{{4.3.5}{20}{Top rows}{lemma.4.3.5}{}}
\@writefile{toc}{\contentsline {subsubsection}{\tocsubsubsection {}{}{Transition steps}}{20}{section*.11}\protected@file@percent }
\newlabel{E:Tn}{{4.3.3}{21}{Transition steps}{equation.4.3.3}{}}
\newlabel{E:P1}{{4.3.4}{21}{Transition steps}{equation.4.3.4}{}}
\newlabel{L:trim}{{4.3.7}{21}{}{lemma.4.3.7}{}}
\newlabel{L:leaves}{{4.3.8}{21}{}{lemma.4.3.8}{}}
\@writefile{toc}{\contentsline {subsubsection}{\tocsubsubsection {}{}{The map $\mathcal  {Z}$}}{21}{section*.12}\protected@file@percent }
\newlabel{definition:zeromap}{{4.3.9}{21}{}{lemma.4.3.9}{}}
\newlabel{lemma:zeroweight}{{4.3.10}{21}{}{lemma.4.3.10}{}}
\citation{GaoHuang2023}
\citation{GaoHuang2023}
\newlabel{lemma:zerotransition}{{4.3.11}{22}{}{lemma.4.3.11}{}}
\newlabel{proposition:zerotrunc}{{4.3.12}{22}{$\zeromap $ erases the last row}{lemma.4.3.12}{}}
\@writefile{lof}{\contentsline {figure}{\numberline {4}{\ignorespaces The bumpless pipe dreams of the RC graph $R=\{(1,2),(1,3),(2,2)\}$ of Figure\nonbreakingspace \ref {fig:zeroing} and of $\mathcal  {Z}(R)$, drawn on a $4\times 4$ grid; blank tiles are shaded and the pipe leaving through each row is labeled at the right. Row $i$ has as many blank tiles as row $i$ of the RC graph has crossings. The two diagrams agree in rows $1$ and $2$ (Proposition\nonbreakingspace \ref {proposition:zerotrunc}); rows $3$ and $4$ of $\Phi (\mathcal  {Z}R)$ contain no blank tile and no $j$-tile, so every pipe there travels straight north and then straight east (Lemma\nonbreakingspace \ref {lemma:bpdcompletion}).}}{22}{figure.caption.13}\protected@file@percent }
\newlabel{fig:bpdzero}{{4}{22}{The bumpless pipe dreams of the RC graph $R=\{(1,2),(1,3),(2,2)\}$ of Figure~\ref {fig:zeroing} and of $\zeromap (R)$, drawn on a $4\times 4$ grid; blank tiles are shaded and the pipe leaving through each row is labeled at the right. Row $i$ has as many blank tiles as row $i$ of the RC graph has crossings. The two diagrams agree in rows $1$ and $2$ (Proposition~\ref {proposition:zerotrunc}); rows $3$ and $4$ of $\Phi (\zeromap R)$ contain no blank tile and no $j$-tile, so every pipe there travels straight north and then straight east (Lemma~\ref {lemma:bpdcompletion})}{figure.caption.13}{}}
\newlabel{lemma:poplocal}{{4.3.13}{23}{Locality of $\nabla $}{lemma.4.3.13}{}}
\newlabel{lemma:ztrim}{{4.3.14}{23}{}{lemma.4.3.14}{}}
\newlabel{theorem:zero_bijection}{{4.3}{23}{}{theorem.4.3}{}}
\newlabel{L:disjoint}{{4.3.15}{23}{}{lemma.4.3.15}{}}
\newlabel{E:first}{{4.3.5}{24}{The map $\zeromap $}{equation.4.3.5}{}}
\newlabel{E:second}{{4.3.6}{24}{The map $\zeromap $}{equation.4.3.6}{}}
\newlabel{L:unique}{{4.3.16}{24}{}{lemma.4.3.16}{}}
\newlabel{lemma:zero_injective}{{4.3.17}{24}{}{lemma.4.3.17}{}}
\newlabel{section:ringproduct}{{4.4}{25}{Definition of the ring product}{subsection.4.4}{}}
\@writefile{toc}{\contentsline {subsection}{\tocsubsection {}{4.4}{Definition of the ring product}}{25}{subsection.4.4}\protected@file@percent }
\newlabel{lemma:lastrow}{{4.4.1}{25}{The last row is determined by its weight}{lemma.4.4.1}{}}
\newlabel{definition:cliptrim}{{4.4.2}{26}{}{lemma.4.4.2}{}}
\newlabel{definition:ringproduct}{{4.4.3}{26}{}{lemma.4.4.3}{}}
\newlabel{lemma:clptrm}{{4.4.4}{26}{Commutativity of $\clp $ and $\trm $}{lemma.4.4.4}{}}
\newlabel{lemma:clpfunc}{{4.4.5}{26}{}{lemma.4.4.5}{}}
\@writefile{toc}{\contentsline {subsection}{\tocsubsection {}{4.5}{Proof of the Littlewood-Richardson rule for the dual Schubert basis}}{28}{subsection.4.5}\protected@file@percent }
\newlabel{lemma:clipfibers}{{4.5.1}{28}{}{lemma.4.5.1}{}}
\newlabel{definition:omega}{{4.5.3}{28}{}{lemma.4.5.3}{}}
\citation{notes}
\newlabel{lemma:schubertrow}{{4.5.5}{29}{}{lemma.4.5.5}{}}
\newlabel{lemma:principalsection}{{4.5.6}{29}{}{lemma.4.5.6}{}}
\newlabel{theorem:semidirect}{{4.5}{29}{}{theorem.4.5}{}}
\newlabel{section:template}{{4.6}{30}{A general Littlewood--Richardson principle}{subsection.4.6}{}}
\@writefile{toc}{\contentsline {subsection}{\tocsubsection {}{4.6}{A general Littlewood--Richardson principle}}{30}{subsection.4.6}\protected@file@percent }
\newlabel{definition:admissible}{{4.6.1}{30}{}{lemma.4.6.1}{}}
\newlabel{theorem:general_lr_template}{{4.6}{30}{General LR principle}{theorem.4.6}{}}
\citation{demazure1974desingularisation}
\citation{demazure1974nouvelle}
\citation{lascoux1990keys}
\citation{reiner1995key}
\citation{lascoux1990keys}
\citation{mason2009explicit}
\citation{kashiwara1993crystal}
\citation{littelmann1995crystal}
\citation{assaf2018demazure}
\citation{morse2016crystal}
\newlabel{remark:general_lr_instances}{{4.6}{31}{}{Item.15}{}}
\@writefile{toc}{\contentsline {section}{\tocsection {}{5}{The ring of RC graph crystals}}{31}{section.5}\protected@file@percent }
\newlabel{section:crystals}{{5}{31}{The ring of RC graph crystals}{section.5}{}}
\citation{polo1989excellent}
\citation{vanderkallen1989longest}
\citation{mathieu1990filtrations}
\citation{joseph1985demazure}
\citation{henrymath}
\citation{morse2016crystal}
\citation{assaf2018demazure}
\citation{assaf2018demazure}
\citation{lascoux1990keys}
\citation{assaf2018demazure}
\citation{hamaker2014relating}
\citation{edelman1987balanced}
\@writefile{toc}{\contentsline {subsection}{\tocsubsection {}{5.1}{Little bumps and the crystal action on RC graphs}}{32}{subsection.5.1}\protected@file@percent }
\newlabel{definition:crystal}{{5.1.1}{32}{Crystal action on RC graphs, highest weights, extremal weights}{lemma.5.1.1}{}}
\newlabel{theorem:keykey}{{5.1}{32}{\cite {lascoux1990keys,assaf2018demazure}}{theorem.5.1}{}}
\newlabel{theorem:littlecrystal}{{5.2}{32}{\cite [Theorem~4.4]{hamaker2014relating}}{theorem.5.2}{}}
\newlabel{corollary:crystaliso}{{5.1.2}{32}{$\zeromap $ intertwines with crystal operators}{lemma.5.1.2}{}}
\citation{edelman1987balanced}
\citation{morse2016crystal}
\citation{hamaker2014relating}
\citation{edelman1987balanced}
\@writefile{lof}{\contentsline {figure}{\numberline {5}{\ignorespaces Edelman--Greene data of an RC graph. The letters of the reading word are inserted from right to left; $P$ is the insertion tableau and $T^-$ is the recording tableau with each entry replaced by the negative of the row of the corresponding crossing. In (a), inserting $3$, then $2$ (which bumps $3$), then $3$ gives $P(R)$ with recording tableau $\begin  {smallmatrix}1&3 2\end  {smallmatrix}$; the crossings inserted first, second and third lie in rows $2,1,1$, so $T^-(R)=\begin  {smallmatrix}-2&-1 -1\end  {smallmatrix}$. The operator $f_1$ moves the crossing $(1,2)$ of $R$ to $(2,3)$ and, via the special bump of $2$ into the row $2\,3$, leaves $P$ unchanged while changing the entry $-1$ in the first row of $T^-$ to $-2$.}}{33}{figure.caption.14}\protected@file@percent }
\newlabel{fig:egexample}{{5}{33}{Edelman--Greene data of an RC graph. The letters of the reading word are inserted from right to left; $P$ is the insertion tableau and $T^-$ is the recording tableau with each entry replaced by the negative of the row of the corresponding crossing. In (a), inserting $3$, then $2$ (which bumps $3$), then $3$ gives $P(R)$ with recording tableau $\begin {smallmatrix}1&3\\2\end {smallmatrix}$; the crossings inserted first, second and third lie in rows $2,1,1$, so $T^-(R)=\begin {smallmatrix}-2&-1\\-1\end {smallmatrix}$. The operator $f_1$ moves the crossing $(1,2)$ of $R$ to $(2,3)$ and, via the special bump of $2$ into the row $2\,3$, leaves $P$ unchanged while changing the entry $-1$ in the first row of $T^-$ to $-2$}{figure.caption.14}{}}
\@writefile{toc}{\contentsline {subsection}{\tocsubsection {}{5.2}{The dual key basis and the ring of RC graph crystals}}{33}{subsection.5.2}\protected@file@percent }
\newlabel{theorem:keyformula}{{5.3}{33}{}{theorem.5.3}{}}
\newlabel{definition:theta}{{5.2.2}{34}{The submodule $\Theta \subseteq \BRC $}{lemma.5.2.2}{}}
\newlabel{lemma:thetaideal}{{5.2.3}{34}{}{lemma.5.2.3}{}}
\newlabel{lemma:keyfibers}{{5.2.6}{34}{Fiber constancy for crystal components}{lemma.5.2.6}{}}
\newlabel{proposition:keyadmissible}{{5.2.7}{35}{}{lemma.5.2.7}{}}
\newlabel{theorem:LRkey}{{5.4}{35}{LR rule for dual keys}{theorem.5.4}{}}
\newlabel{example:LRkey}{{5.2.8}{35}{Computing $\kappa _a^*\sqcupplus \kappa _b^*$ for $a=(0,1,2)$, $b=(0,1,0,2)$}{lemma.5.2.8}{}}
\@writefile{lof}{\contentsline {figure}{\numberline {6}{\ignorespaces The two $7$-row RC graphs $R_1$ (left) and $R_2$ (right) in $\mathrm  {Dem}_c$ satisfying $\mathtt  {clip}^3(R_i)\in \mathrm  {Dem}_a$ and $\mathtt  {trim}^3(R_i)\in \mathrm  {Dem}_b$. The outlined box marks the top $3$ rows (the clip region).}}{35}{figure.caption.15}\protected@file@percent }
\newlabel{fig:LRkey_example}{{6}{35}{The two $7$-row RC graphs $R_1$ (left) and $R_2$ (right) in $\mathrm {Dem}_c$ satisfying $\clp ^3(R_i)\in \mathrm {Dem}_a$ and $\trm ^3(R_i)\in \mathrm {Dem}_b$. The outlined box marks the top $3$ rows (the clip region)}{figure.caption.15}{}}
\citation{nadeau2024forest}
\@writefile{toc}{\contentsline {section}{\tocsection {}{6}{The ring of forest classes of RC graphs}}{36}{section.6}\protected@file@percent }
\newlabel{section:forest}{{6}{36}{The ring of forest classes of RC graphs}{section.6}{}}
\@writefile{toc}{\contentsline {subsection}{\tocsubsection {}{6.1}{Indexed forests, forest invariant}}{36}{subsection.6.1}\protected@file@percent }
\citation{nadeau2024ppartitions}
\citation{nadeau2024forest}
\citation{nadeau2024forest}
\citation{nadeau2024forest}
\citation{nadeau2024forest}
\newlabel{definition:forestequiv}{{6.1.3}{37}{Injectification, $\equiv _\forest $}{lemma.6.1.3}{}}
\newlabel{theorem:forestforest}{{6.2}{37}{\cite [Theorem~1.4]{nadeau2024forest}}{theorem.6.2}{}}
\citation{bergeron2025qsymflag}
\citation{nadeau2024quasisymmetric}
\citation{nadeau2024quasisymmetric}
\citation{nadeau2024quasisymmetric}
\citation{nadeau2024quasisymmetric}
\citation{nadeau2024quasisymmetric}
\citation{nadeau2024quasisymmetric}
\citation{postnikov2009chains}
\citation{bergeron2025qsymflag}
\citation{bergeron2025qsymflag}
\citation{nadeau2024quasisymmetric}
\citation{nadeau2024quasisymmetric}
\citation{nadeau2024quasisymmetric}
\citation{nadeau2024quasisymmetric}
\citation{bergeron2025qsymflag}
\@writefile{toc}{\contentsline {subsection}{\tocsubsection {}{6.2}{The dual forest basis and the ring of forest classes of RC graphs}}{38}{subsection.6.2}\protected@file@percent }
\newlabel{proposition:forestdualcoefficients}{{6.2.2}{38}{}{lemma.6.2.2}{}}
\newlabel{theorem:forestformula}{{6.3}{38}{}{theorem.6.3}{}}
\citation{nadeau2024forest}
\citation{nadeau2024forest}
\citation{guo2026forest}
\citation{nadeau2024forest}
\citation{nadeau2024forest}
\newlabel{definition:gamma}{{6.2.4}{39}{The submodule $\Gamma \subseteq \BRC $}{lemma.6.2.4}{}}
\newlabel{lemma:forestconcat}{{6.2.5}{39}{$\equiv _\forest $ and concatenation}{lemma.6.2.5}{}}
\newlabel{lemma:forestzeropop}{{6.2.6}{39}{Deleting the first crossing commutes with the zero map}{lemma.6.2.6}{}}
\newlabel{lemma:forestzeroleadingstep}{{6.2.7}{40}{Leading-letter transport in a reversed transition}{lemma.6.2.7}{}}
\newlabel{corollary:forestzeroleadingunique}{{6.2.8}{40}{The lifted leading letter is unique}{lemma.6.2.8}{}}
\citation{nadeau2024forest}
\newlabel{lemma:foresthighrows}{{6.2.9}{41}{A high-supported forest has a row-supported representative}{lemma.6.2.9}{}}
\newlabel{lemma:forestseparatedshuffle}{{6.2.10}{41}{Separated-support shuffles}{lemma.6.2.10}{}}
\newlabel{proposition:forestcompatiblechain}{{6.2.11}{41}{Fixed-chain lifts with a compatible right factor}{lemma.6.2.11}{}}
\newlabel{lemma:forestfibers}{{6.2.12}{43}{Fiber constancy for forest classes}{lemma.6.2.12}{}}
\newlabel{lemma:forestideal}{{6.2.13}{43}{}{lemma.6.2.13}{}}
\citation{nadeau2024forest}
\citation{nadeau2024forest}
\citation{bergeron2025qsymflag}
\citation{bergeron2025qsymflag}
\citation{bergeron2025qsymflag}
\citation{bergeron2025qsymflag}
\citation{nadeau2024forest}
\newlabel{proposition:forestadmissible}{{6.2.16}{44}{}{lemma.6.2.16}{}}
\newlabel{theorem:LRforest}{{6.4}{44}{LR rule for dual forests}{theorem.6.4}{}}
\newlabel{example:LRforest}{{6.2.17}{44}{Computing $\forest _a^*\sqcupplus \forest _b^*$ for $a=(0,0,2,0,2)$, $b=(0,1,0,0,2)$}{lemma.6.2.17}{}}
\@writefile{lof}{\contentsline {figure}{\numberline {7}{\ignorespaces The two $10$-row witness pipe dreams in the forest class for $c$.}}{45}{figure.caption.16}\protected@file@percent }
\newlabel{fig:forest_lr_witnesses}{{7}{45}{The two $10$-row witness pipe dreams in the forest class for $c$}{figure.caption.16}{}}
\@writefile{lof}{\contentsline {figure}{\numberline {8}{\ignorespaces Common clip/trim pair for both witnesses $R_1,R_2$.}}{45}{figure.caption.17}\protected@file@percent }
\newlabel{fig:forest_lr_cliptrim}{{8}{45}{Common clip/trim pair for both witnesses $R_1,R_2$}{figure.caption.17}{}}
\newlabel{subsection:qflbranch}{{6.3}{45}{Geometric interpretation: branching on the quasisymmetric flag variety}{subsection.6.3}{}}
\@writefile{toc}{\contentsline {subsection}{\tocsubsection {}{6.3}{Geometric interpretation: branching on the quasisymmetric flag variety}}{45}{subsection.6.3}\protected@file@percent }
\citation{ressayre2009branching}
\citation{assaf2017slide}
\citation{assaf2017slide}
\citation{nadeau2024ppartitions}
\@writefile{toc}{\contentsline {section}{\tocsection {}{7}{Product rule for dual slide polynomials}}{46}{section.7}\protected@file@percent }
\newlabel{section:slide}{{7}{46}{Product rule for dual slide polynomials}{section.7}{}}
\@writefile{toc}{\contentsline {subsection}{\tocsubsection {}{7.1}{Slide polynomials}}{46}{subsection.7.1}\protected@file@percent }
\newlabel{lemma:slideideal}{{7.1.4}{46}{}{lemma.7.1.4}{}}
\citation{assaf2017slide}
\citation{assaf2017slide}
\newlabel{lemma:slidefibers}{{7.1.7}{47}{Fiber constancy for word classes}{lemma.7.1.7}{}}
\newlabel{proposition:slidebasis}{{7.1.8}{47}{}{lemma.7.1.8}{}}
\citation{nadeau2024forest}
\citation{bergeron2025qsymflag}
\newlabel{theorem:lrfslide}{{7.2}{48}{}{theorem.7.2}{}}
\@writefile{toc}{\contentsline {section}{\tocsection {}{}{Notation and conventions}}{48}{section*.18}\protected@file@percent }
\newlabel{section:notation}{{7.1}{48}{Notation and conventions}{section*.18}{}}
\@writefile{toc}{\contentsline {subsubsection}{\tocsubsubsection {}{}{Permutations and combinatorial objects}}{48}{section*.19}\protected@file@percent }
\@writefile{toc}{\contentsline {subsubsection}{\tocsubsubsection {}{}{Row maps on RC graphs}}{48}{section*.20}\protected@file@percent }
\@writefile{toc}{\contentsline {subsubsection}{\tocsubsubsection {}{}{Crystal and forest structures}}{48}{section*.21}\protected@file@percent }
\@writefile{toc}{\contentsline {subsubsection}{\tocsubsubsection {}{}{The Schubert bialgebra and its dual}}{49}{section*.22}\protected@file@percent }
\@writefile{toc}{\contentsline {subsubsection}{\tocsubsubsection {}{}{Structure constants}}{49}{section*.23}\protected@file@percent }
\@writefile{toc}{\contentsline {subsubsection}{\tocsubsubsection {}{}{The ring of RC graphs and its quotients}}{49}{section*.24}\protected@file@percent }
\citation{bergeron2025qsymflag}
\citation{bergeron2025qsymflag}
\citation{bergeron2025qsymflag}
\citation{nadeau2024quasisymmetric}
\citation{nadeau2024quasisymmetric}
\citation{samuelmolev}
\@writefile{toc}{\contentsline {subsubsection}{\tocsubsubsection {}{}{Slide polynomials and the general LR template}}{50}{section*.25}\protected@file@percent }
\@writefile{toc}{\contentsline {subsubsection}{\tocsubsubsection {}{}{Quasisymmetric flag variety (\S  \ref {subsection:qflbranch})}}{50}{section*.26}\protected@file@percent }
\@writefile{toc}{\contentsline {section}{\tocsection {Appendix}{A}{Proof of Proposition \ref {proposition:pullindex}}}{50}{appendix.A}\protected@file@percent }
\newlabel{proposition:pieri}{{A.1}{50}{Special case of Pieri formula {\cite [Theorem~7.1]{samuelmolev}}}{proposition2.A.1}{}}
\bibstyle{acm}
\bibdata{schubert_bialgebra,schubert_bialgebra_bpd}
\bibcite{assaf2018demazure}{1}
\bibcite{assaf2017slide}{2}
\bibcite{bbrc}{3}
\bibcite{bergeron2025qsymflag}{4}
\bibcite{coprod}{5}
\bibcite{bsskew}{6}
\bibcite{billey2019bijective}{7}
\bibcite{demazure1974desingularisation}{8}
\bibcite{demazure1974nouvelle}{9}
\bibcite{edelman1987balanced}{10}
\bibcite{GaoHuang2023}{11}
\bibcite{guo2026forest}{12}
\bibcite{hamaker2014relating}{13}
\bibcite{henrymath}{14}
\bibcite{huang2020monk}{15}
\bibcite{HuangShimozonoYu2024}{16}
\bibcite{joseph1985demazure}{17}
\bibcite{kashiwara1993crystal}{18}
\bibcite{lam2021back}{19}
\bibcite{lascoux1990keys}{20}
\bibcite{littelmann1995crystal}{21}
\bibcite{little2003combinatorial}{22}
\bibcite{notes}{23}
\bibcite{mason2009explicit}{24}
\bibcite{mathieu1990filtrations}{25}
\bibcite{morse2016crystal}{26}
\bibcite{nadeau2024quasisymmetric}{27}
\bibcite{nadeau2024forest}{28}
\bibcite{nadeau2024ppartitions}{29}
\bibcite{polo1989excellent}{30}
\bibcite{postnikov2009chains}{31}
\bibcite{reiner1995key}{32}
\bibcite{ressayre2009branching}{33}
\@writefile{toc}{\contentsline {section}{\tocsection {Appendix}{}{Acknowledgements}}{52}{section*.27}\protected@file@percent }
\@writefile{toc}{\contentsline {section}{\tocsection {Appendix}{}{References}}{52}{section*.28}\protected@file@percent }
\bibcite{samuelleibniz}{34}
\bibcite{samuelmolev}{35}
\bibcite{samuelforest}{36}
\bibcite{schubmult}{37}
\bibcite{sottile}{38}
\bibcite{vanderkallen1989longest}{39}
\bibcite{Weigandt2021}{40}
\bibcite{Yu2025EmbeddingBP}{41}
\newlabel{tocindent-1}{0pt}
\newlabel{tocindent0}{58.61125pt}
\newlabel{tocindent1}{66.11127pt}
\newlabel{tocindent2}{29.38896pt}
\newlabel{tocindent3}{0pt}
\gdef \@abspage@last{53}
\end{filecontents*}
